\documentclass[a4paper,12pt]{article}
\usepackage[T1]{fontenc}
\usepackage[brazil]{babel}
\usepackage{amsmath,amssymb}
\usepackage{graphicx}
\usepackage{booktabs}
\usepackage[colorlinks=true,allcolors=blue]{hyperref}

\title{{{TITLE}}}
\author{Seu Nome}
\date{\today}

\begin{document}

\maketitle

\begin{abstract}
Este é um projeto de exemplo. Edite o texto à esquerda e veja o PDF à direita.
\end{abstract}

\section{Introdução}\label{sec:intro}
O \LaTeX{} cuida da numeração, das referências cruzadas e da bibliografia.
Veja a Figura~\ref{fig:exemplo}, a Tabela~\ref{tab:exemplo} e a
Equação~\eqref{eq:euler}. Uma citação fica assim~\cite{knuth1984}.

\section{Figuras}
\begin{figure}[ht]
  \centering
  \includegraphics[width=0.5\linewidth]{example-image}
  \caption{Uma imagem de exemplo.}
  \label{fig:exemplo}
\end{figure}

\section{Tabelas}
\begin{table}[ht]
  \centering
  \begin{tabular}{lrr}
    \toprule
    Item & Quantidade & Preço \\
    \midrule
    Café & 2 & 7,50 \\
    Pão  & 6 & 4,20 \\
    \bottomrule
  \end{tabular}
  \caption{Uma tabela simples.}
  \label{tab:exemplo}
\end{table}

\section{Matemática}
Fórmulas no texto, como $a^2 + b^2 = c^2$, ou em destaque:
\begin{equation}
  e^{i\pi} + 1 = 0 \label{eq:euler}
\end{equation}

\bibliographystyle{plain}
\bibliography{referencias}

\end{document}
